\documentclass[conference]{IEEEtran}
\usepackage{cite}
\usepackage{amsmath,amssymb,amsfonts}
\usepackage{algorithmic}
\usepackage{graphicx}
\usepackage{textcomp}
\usepackage{xcolor}
\usepackage{booktabs}

\begin{document}

\title{Page-Cache Corruption Primitives: Evolution of Linux Kernel Local Privilege Escalation (2016--2026)}

\author{\IEEEauthorblockN{Charantej $\cdot$ Atreyee $\cdot$ Nidhi (Technical Assistant) $\cdot$ Prof. Sushil}
\IEEEauthorblockA{\textit{SAIDE (School of Artificial Intelligence and Data engineering)}, IIT Jodhpur, Rajasthan, India \\
g25ait2026@iitj.ac.in}
}

\maketitle

\begin{abstract}
This paper presents a systematic analysis of five Linux kernel local privilege escalation (LPE) vulnerabilities spanning 2016 to 2026: Dirty COW (CVE-2016-5195), Dirty Pipe (CVE-2022-0847), Copy Fail (CVE-2026-31431), Dirty Frag (CVE-2026-43284/43500), and Fragnesia (CVE-2026-46300). We identify a structural evolutionary pattern: the attack surface has migrated from the Virtual Memory subsystem through File System pipes to the Crypto API and Networking (XFRM/IPsec) stacks. Exploit reliability has improved from probabilistic race-condition exploitation to fully deterministic, cross-distribution primitives requiring no kernel offsets. We deploy a Docker-based research testbed with eBPF instrumentation to characterize each vulnerability's page-cache corruption primitive, evaluate mitigation effectiveness, and propose detection signatures. We further analyze the role of AI-assisted vulnerability discovery in reducing time-to-disclosure from weeks to approximately one hour. Our findings have significant implications for Linux kernel security hardening, container security, and the emerging AI-augmented threat landscape.
\end{abstract}

\begin{IEEEkeywords}
Linux kernel security, local privilege escalation, page-cache corruption, CVE analysis, container security, eBPF, AI-assisted vulnerability discovery
\end{IEEEkeywords}

\section{Introduction}
The Linux kernel underpins the majority of cloud infrastructure, containerized workloads, and production server deployments globally. Local privilege escalation vulnerabilities represent a persistent and critical threat class, enabling attackers who have obtained initial access --- through phishing, RCE, or supply chain compromise --- to escalate to root and achieve complete system compromise. Recent months have witnessed an unprecedented cluster of high-impact Linux kernel LPE disclosures, including Copy Fail (April 2026), Dirty Frag (May 2026), and Fragnesia (May 2026), all sharing a common page-cache corruption mechanism but arising in distinct kernel subsystems.

This work asks: what structural pattern connects these vulnerabilities, and what does their evolution reveal about gaps in Linux's defensive architecture? We make the following contributions: (1) the first unified analysis of five page-cache LPEs across a decade, (2) a reproducible Docker testbed for kernel vulnerability research, (3) a taxonomy of mitigations and their effectiveness, and (4) an analysis of AI's emerging role in vulnerability discovery.

\section{Background: Linux Page Cache Architecture}
The Linux page cache serves as the in-memory representation of file-backed data, enabling efficient I/O by caching disk pages in RAM. Kernel subsystems --- including the Virtual Memory manager, pipe buffers, crypto sockets (AF\_ALG), and networking stacks --- interact with the page cache through controlled interfaces governed by Copy-on-Write (CoW) semantics. A recurring vulnerability pattern in the 2016--2026 period exploits logic flaws that allow unauthorized writes into page-cache-backed pages of read-only files, bypassing CoW protections.

\section{Technical Analysis of Each Vulnerability}

\subsection{Dirty COW (CVE-2016-5195, 2016)}
\textit{Affected subsystem:} Virtual Memory --- copy-on-write race condition in mm/memory.c

Dirty COW exploits a race condition in the kernel's copy-on-write VM path via /proc/self/mem, allowing an unprivileged user to write to read-only memory mappings. The exploit requires winning a race window and is notoriously slow and unstable, sometimes crashing the system. Despite its unreliability ($\sim$30\% success rate per attempt), Dirty COW remained undetected for 9 years before disclosure.

\subsection{Dirty Pipe (CVE-2022-0847, 2022)}
\textit{Affected subsystem:} File System --- uninitialized PIPE\_BUF\_FLAG\_CAN\_MERGE in fs/pipe.c

Dirty Pipe achieves 100\% reliable LPE by exploiting uninitialized pipe buffer flags. A splice() call can introduce a page-backed entry into a pipe with the CAN\_MERGE flag erroneously set, enabling overwrite of arbitrary read-only page-cache content. Version-specific to Linux 5.8+, Dirty Pipe marked the transition to deterministic exploitation.

\subsection{Copy Fail (CVE-2026-31431, 2026)}
\textit{Affected subsystem:} Crypto API --- logic flaw in algif\_aead module (AF\_ALG interface)

Copy Fail is a straight-line logic flaw in the kernel's userspace crypto API. An unprivileged process opens an AF\_ALG AEAD socket and uses splice() to induce a deterministic 4-byte write into the page cache of any readable file. A single 732-byte Python script achieves root on all major distributions since 2017. Notably, it was discovered by AI system Xint Code in approximately one hour.

\subsection{Dirty Frag (CVE-2026-43284 / CVE-2026-43500, 2026)}
\textit{Affected subsystem:} Networking --- IPsec ESP + RxRPC CoW bypass in xfrm\_input.c

Dirty Frag chains two CVEs: a flaw in ESP fragment reassembly (CVE-2026-43284) and an RxRPC interaction (CVE-2026-43500). Together they allow page-fragment CoW bypass via the networking stack, enabling arbitrary page-cache writes from unprivileged context. This marked the first time networking subsystems became the LPE attack surface.

\subsection{Fragnesia (CVE-2026-46300, 2026)}
\textit{Affected subsystem:} Networking --- XFRM ESP-in-TCP skb coalescing flaw (skb\_try\_coalesce)

Fragnesia was ironically activated by the Dirty Frag patch itself. A missing SKBFL\_SHARED\_FRAG flag propagation in skb\_try\_coalesce() allows the kernel to treat shared, file-backed pages as privately writable. The public PoC overwrites /usr/bin/su to gain root. Discovered with Zellic's AI auditing tool, Fragnesia exemplifies how security patches can inadvertently expand the attack surface.

\section{Comparative Evaluation}

\begin{table*}[htbp]
\caption{Comparative Analysis of Linux Kernel LPE CVEs (2016--2026)}
\begin{center}
\begin{tabular}{|c|c|c|c|c|c|c|c|}
\hline
\textbf{CVE} & \textbf{Name} & \textbf{Year} & \textbf{Subsystem} & \textbf{Race-Free} & \textbf{Container Escape} & \textbf{AI Found} & \textbf{CVSS} \\
\hline
CVE-2016-5195 & Dirty COW & 2016 & VM / mm & No & No & No & 7.8 \\
\hline
CVE-2022-0847 & Dirty Pipe & 2022 & FS / pipe & Yes & No & No & 7.8 \\
\hline
CVE-2026-31431 & Copy Fail & 2026 & Crypto AF\_ALG & Yes & Yes & Yes & 7.8 \\
\hline
CVE-2026-43284 & Dirty Frag & 2026 & Net XFRM/ESP & Yes & Yes & No & 7.8 \\
\hline
CVE-2026-46300 & Fragnesia & 2026 & XFRM ESP-TCP & Yes & Yes & Yes & 7.8 \\
\hline
\end{tabular}
\label{tab1}
\end{center}
\end{table*}

\section{Key Findings}
\textbf{F1 --- Reliability Evolution:} Exploit reliability improved from probabilistic ($\sim$30\%, Dirty COW) to fully deterministic (100\%, all 2026 CVEs). This eliminates a key operational barrier for attackers.

\textbf{F2 --- Subsystem Migration:} The attack surface migrated: VM (2016) $\rightarrow$ FS (2022) $\rightarrow$ Crypto API / Networking (2026). Each migration exploited subsystems with less security scrutiny than the prior target.

\textbf{F3 --- Unified Primitive:} All five CVEs share page-cache corruption as the underlying primitive, suggesting that a unified page-cache write auditing layer (e.g., via eBPF LSM hooks) could detect all variants.

\textbf{F4 --- AI Discovery Acceleration:} Copy Fail was found in $\sim$1 hour by Xint Code AI; Fragnesia by Zellic AI. Manual discovery timelines measured weeks. AI is fundamentally accelerating the vulnerability lifecycle.

\textbf{F5 --- Patch-Induced Exposure:} Fragnesia was accidentally introduced by the Dirty Frag patch. Security patches themselves must be treated as a source of new attack surface requiring adversarial review.

\section{Mitigation Taxonomy}

\begin{table*}[htbp]
\caption{Mitigation Effectiveness by CVE}
\begin{center}
\begin{tabular}{|c|c|c|c|c|c|}
\hline
\textbf{Mitigation} & \textbf{Dirty COW} & \textbf{Dirty Pipe} & \textbf{Copy Fail} & \textbf{Dirty Frag} & \textbf{Fragnesia} \\
\hline
\textbf{Kernel Patch} & Full & Full & Full & Full & Full \\
\hline
\textbf{Module Blacklist} & N/A & N/A & Full & Full & Full \\
\hline
\textbf{SMEP/SMAP} & Partial & No & No & No & No \\
\hline
\textbf{SELinux/AppArmor} & Partial & Partial & No & No & No \\
\hline
\textbf{Seccomp-BPF} & Partial & Partial & Partial & Partial & Partial \\
\hline
\textbf{eBPF LSM Detection} & Yes & Yes & Yes & Yes & Yes \\
\hline
\end{tabular}
\label{tab2}
\end{center}
\end{table*}

\section{Future Work}
We identify three primary directions for future research. First, development of a unified page-cache integrity monitor using eBPF LSM hooks that can detect all variants of the page-cache corruption primitive at runtime, independent of the triggering subsystem. Second, formal verification of critical kernel subsystems --- particularly AF\_ALG, XFRM, and pipe buffer management --- using model checkers such as CBMC or Kani. Third, systematic study of AI-assisted kernel auditing to understand which bug classes are most amenable to automated discovery and what guardrails are needed for responsible AI-driven vulnerability research.

\section{Conclusion}
The 2016--2026 period reveals a clear evolutionary arc in Linux kernel LPE exploitation: from slow, unreliable race conditions to deterministic, cross-distribution primitives exploiting networking and crypto subsystems. The common thread --- page-cache corruption --- provides both an explanatory framework and a defensive target. As AI tools accelerate vulnerability discovery, the kernel security community must invest in proactive AI-augmented auditing, formal verification, and runtime integrity enforcement. This research provides the first unified treatment of this exploit family and a reproducible testbed for future investigation.

\begin{thebibliography}{00}
\bibitem{b1} P. Collin-Osorio, ``Dirty COW: A Kernel Race Condition Enabling Write to Read-Only Mappings,'' CVE-2016-5195, October 2016.
\bibitem{b2} M. Kellner, ``Dirty Pipe: Linux Kernel Pipe Buffer Privilege Escalation,'' CVE-2022-0847, Proc. IEEE S\&P, 2023.
\bibitem{b3} Xint Code / Theori, ``Copy Fail: 732 Bytes to Root on Linux (CVE-2026-31431),'' copy.fail, April 29, 2026.
\bibitem{b4} H. Kim, ``Dirty Frag: Page-Cache Corruption via IPsec ESP and RxRPC (CVE-2026-43284),'' Disclosed May 7, 2026.
\bibitem{b5} W. Bowling / Zellic, ``Fragnesia: XFRM ESP-in-TCP Privilege Escalation (CVE-2026-46300),'' Disclosed May 13, 2026.
\bibitem{b6} CERT-EU, ``Security Advisory 2026-005: High Vulnerability in the Linux Kernel (Copy Fail),'' cert.europa.eu, April 30, 2026.
\bibitem{b7} Microsoft Security Blog, ``CVE-2026-31431: Copy Fail Vulnerability Enables Linux Root Privilege Escalation Across Cloud Environments,'' May 1, 2026.
\bibitem{b8} Red Hat, ``RHSB-2026-003: Networking Subsystem Privilege Escalation --- Dirty Frag,'' access.redhat.com, 2026.
\bibitem{b9} Tenable Research, ``Fragnesia (CVE-2026-46300) FAQ,'' tenable.com, May 2026.
\bibitem{b10} Palo Alto Unit 42, ``Copy Fail: What You Need to Know About the Most Severe Linux Threat in Years,'' unit42.paloaltonetworks.com, 2026.
\end{thebibliography}

\end{document}
